\begin{mistake}{2026-05-28}{05}

\begin{problemcontent}
设 \(P(x),Q(x)\) 在 \((-
\infty,+\infty)\) 内连续，且均以 \(T\) 为周期，函数 \(y=y(x)\) 是微分方程
\[
\frac{dy}{dx}+P(x)y=Q(x)
\]
的解。判断条件 \(y(0)=y(T)\) 是 \(y=y(x)\) 以 \(T\) 为周期的什么条件？
\[
\begin{array}{ll}
\text{(A) 必要非充分条件} & \text{(B) 充分非必要条件}\\
\text{(C) 充要条件} & \text{(D) 既非充分又非必要条件}
\end{array}
\]
\end{problemcontent}

\begin{solutioncontent}
答案为 \(\text{(C)}\)。

必要性：若 \(y(x)\) 以 \(T\) 为周期，则对任意 \(x\) 有 \(y(x+T)=y(x)\)。令 \(x=0\)，得
\[
y(T)=y(0),
\]
所以 \(y(0)=y(T)\) 是必要条件。

充分性：设 \(y(0)=y(T)\)。构造
\[
Y(x)=y(x+T).
\]
因为 \(P(x),Q(x)\) 均以 \(T\) 为周期，所以 \(P(x+T)=P(x)\)，\(Q(x+T)=Q(x)\)。又因 \(y(x)\) 是原方程的解，故
\[
y'(x+T)+P(x+T)y(x+T)=Q(x+T).
\]
于是
\[
\begin{aligned}
Y'(x)+P(x)Y(x)
&=y'(x+T)+P(x)y(x+T)\\
&=y'(x+T)+P(x+T)y(x+T)\\
&=Q(x+T)=Q(x).
\end{aligned}
\]
因此 \(Y(x)\) 也是微分方程 \(y'+P(x)y=Q(x)\) 的解。

在 \(x=0\) 处，
\[
Y(0)=y(T)=y(0).
\]
由于 \(P(x),Q(x)\) 在 \((-
\infty,+\infty)\) 内连续，一阶线性微分方程初值问题解唯一，所以
\[
Y(x)\equiv y(x),
\]
即
\[
y(x+T)=y(x).
\]
故 \(y(x)\) 以 \(T\) 为周期。因此 \(y(0)=y(T)\) 是充要条件。
\end{solutioncontent}

\mistakeknowledge{
\begin{itemize}
  \item \textbf{核心公式/定理}：一阶线性微分方程 \(y'+P(x)y=Q(x)\) 中，若 \(P,Q\) 连续，则给定初值 \(y(x_0)=y_0\) 后解唯一。
  \item \textbf{易错点}：充分性不能只由 \(y(0)=y(T)\) 直接断言周期性，必须构造 \(Y(x)=y(x+T)\)，证明其与 \(y(x)\) 满足同一方程和同一初值。
  \item \textbf{关联知识}：周期系数问题常用“平移构造解 + 唯一性定理”；若 \(P,Q\) 以 \(T\) 为周期，则 \(P(x+T)=P(x)\)，\(Q(x+T)=Q(x)\)。
\end{itemize}}

\end{mistake}
